\chapter{Discussion, Product Context, and Conclusion}

\section{Error Analysis Plan}
The error analysis should explain why query rewriting helps or fails. Possible failure classes include:

\begin{itemize}
    \item \textbf{Missing source}: the correct official source is not in the indexed corpus;
    \item \textbf{Bad source granularity}: the document is indexed, but the relevant information is hidden in a long segment or split across segments;
    \item \textbf{Ambiguous query}: the question is course-specific but does not specify the degree programme;
    \item \textbf{Rewrite drift}: the LLM changes the original intent or adds assumptions;
    \item \textbf{Lexical mismatch remains}: the rewritten query still does not match the official wording;
    \item \textbf{Duplicate yearly documents}: old and new versions compete in the ranking;
    \item \textbf{Right retrieval, wrong downstream answer}: relevant mainly for the product prototype, not for the core retrieval experiment.
\end{itemize}

\section{Product Context}
This thesis is extracted from a broader product-oriented prototype for a university bureaucracy assistant. The broader system includes controlled source acquisition, manifest/versioning, HTML/PDF extraction, BM25 indexing, a backend service, a frontend interface, and debugging tools for inspecting acquisition and extraction.

The product-oriented goal is to support students and administrative staff by providing a trustworthy interface over official university sources. However, the thesis isolates one measurable component: whether LLM query rewriting improves retrieval quality.

\section{Limitations}
The current thesis focuses on retrieval rather than full answer generation. This keeps the experiment measurable, but it does not fully evaluate chatbot behavior. Citation correctness, refusal behavior, answer faithfulness, and user satisfaction are therefore outside the core evaluation unless added as secondary analysis.

Another limitation is source coverage. If the correct document is missing from the corpus, query rewriting cannot solve the problem. In such cases, the limitation belongs to acquisition and corpus completeness, not only to retrieval.

\section{Future Work}
Future work may include:

\begin{itemize}
    \item hybrid retrieval combining BM25 and vector search;
    \item reranking retrieved results;
    \item structured databases for hard constraints such as ECTS requirements;
    \item stronger citation and refusal evaluation for generated answers;
    \item broader source acquisition across additional university domains;
    \item user studies with students and administrative staff;
    \item full assistant evaluation including answer correctness and faithfulness.
\end{itemize}

\section{Conclusion}
This thesis evaluates whether a low-cost LLM can improve deterministic retrieval by rewriting informal student queries into better BM25 search queries. The broader motivation is that, in institutional question answering, answer fluency is not sufficient: correctness depends on retrieving the right official context.

The central contribution is therefore not a generic chatbot, but an empirical study of a lightweight AI layer that may improve a hard-coded retrieval pipeline while preserving inspectability and provenance.
\par
